\clearpage\section*{Correction EV3-A}\label{EV3-A-CORR}
\textbf{QCM :} 1 A, 2 B, 3 A, 4 C, 5 C, 6 D, 7 B, 8 D, 9 D, 10 B, 11 A, 12 B. Chaque réponse correcte : 0,5 point ; total 6.
\paragraph{Vérifications des automatismes.} Priorités ; somme de fractions ; soustraction des exposants ; coefficient multiplicateur 1,2 ; division par 0,8 ; produit 1,25 fois 0,8 ; résolution affine ; accroissement vertical divisé par l’horizontal ; moyenne ; angle divisé par 360 ; complément à 1 ; substitution entre parenthèses.
\subsection*{Exercice 1 -- Loterie et dispersion}
\begin{enumerate}
\item X vaut -3,1,9 avec les mêmes probabilités. \textit{Barème : 1 point. Résultat et justification élémentaire attendus ; demi-point possible pour démarche correcte inachevée.}
\item $E(X)=-\dfrac{3}{2}+\dfrac{3}{8}+\dfrac{9}{8}=0$. $E(X^2)=\dfrac{9}{2}+\dfrac{3}{8}+\dfrac{81}{8}=15$ ; variance 15. \textit{Barème : 2 points. Méthode ou justification : 1 ; calcul et conclusion : 1.}
\item Gain net 2X, espérance 0 ; écart type $2\sqrt{15}$. \textit{Barème : 2 points. Méthode ou justification : 1 ; calcul et conclusion : 1.}
\end{enumerate}
\subsection*{Exercice 2 -- Dérivée et tangente horizontale}
\begin{enumerate}
\item $f(0)=-3$ ; $f\prime=(2x-1)e^x$, $f\prime(0)=-1$. \textit{Barème : 1 point. Résultat et justification élémentaire attendus ; demi-point possible pour démarche correcte inachevée.}
\item Décroît jusqu’à $\dfrac{1}{2}$, croît ensuite ; minimum $-2e^{\dfrac{1}{2}}$. \textit{Barème : 2 points. Méthode ou justification : 1 ; calcul et conclusion : 1.}
\item $y=-x-3$ ; contact (0;-3), sous l’axe horizontal. \textit{Barème : 2 points. Méthode ou justification : 1 ; calcul et conclusion : 1.}
\end{enumerate}
\subsection*{Exercice 3 -- Angle et longueur}
\begin{enumerate}
\item AB=4, $AC=\sqrt{4+12}=4$. \textit{Barème : 1 point. Résultat et justification élémentaire attendus ; demi-point possible pour démarche correcte inachevée.}
\item Produit 8 ; cosinus $\dfrac{8}{16}=\dfrac{1}{2}$, angle $\dfrac{\pi}{3}$. \textit{Barème : 1 point. Résultat et justification élémentaire attendus ; demi-point possible pour démarche correcte inachevée.}
\item $BC=\sqrt{4+12}=4$ : triangle équilatéral. \textit{Barème : 2 points. Méthode ou justification : 1 ; calcul et conclusion : 1.}
\end{enumerate}
\paragraph{Étalonnage.} Durée cible : 20–25 min pour les automatismes, 80–85 min pour les exercices, 10–15 min de relecture. Ces durées sont des estimations éditoriales à mesurer en classe. Accepter toute méthode correcte ; ne pas sanctionner deux fois une erreur numérique propagée avec cohérence. Le barème proposé ne vaut pas décision d’harmonisation de l’équipe.
